\begin{python0}
from solutions import *; clear()
\end{python0}
\sectionthree{Program 6. Catching Multiple Exceptions}

Of course there's no reason to believe that in a chunk
of code (say in a function) there is only one possible error. The
following show you that a function, in the example that would be
\texttt{main()}, can be a catcher to two different types of exceptions.



\begin{ex}
Recall a previous exercise:

\begin{console}[commandchars=\~\@\$]
#include <iostream>
#include <cmath>
#include <ctime>

~EMPHASIZE@class StackException$
~EMPHASIZE@{};$

class Stack
{     
public:
      ...
      void push(int i)
      {    
           ~EMPHASIZE@if (size == 5)$
                 ~EMPHASIZE@throw StackException();$
           ...
      }
      int pop()
      {   
          ~EMPHASIZE@if (size == 0)$
                 ~EMPHASIZE@throw StackException();$
          ...
      }
  ...
};



int main()
{   
    srand((unsigned) time(NULL));
    Stack stack;

    try
    {   
        while (1)
        {     
              int option;
              std::cin >> option;
              switch (option)
              {
                    case 0: stack.push(rand()); break;
                    case 1: stack.pop(); break;}}}}
              }
              std::cout << stack << std::endl;
         }
     }
     catch (~EMPHASIZE@StackException e$)
     {
          std::cout << "stack error" << std::endl;
     }
} 
\end{console}

Modify the program by replacing StackException with two new classes,
\texttt{StackOverflowException} and \texttt{StackUnderflowException}, to
handle the two different errors. The code in \texttt{main()} should catch
both errors.
\end{ex}
By the way it's a good idea to keep exceptions which are
thrown only in a class together with the class itself in the same header
file. It makes the exception class easier to find.

